% !TEX root = ../main.tex
\chapter{Introducció}
\label{chap:introduccio}

\section{Contextualització i motivació}

Pokémon és un joc d'estratègia per torns caracteritzat per la informació parcial, les accions simultànies i les transicions estocàstiques. A cada torn, els jugadors han de decidir entre executar un moviment o substituir el Pokémon actiu sense conèixer amb certesa la configuració de l'equip rival, mentre que factors com la precisió dels atacs, els cops crítics i els efectes secundaris introdueixen aleatorietat inherent al combat.

En aquest context, un \textbf{agent} és un sistema autònom que percep l'estat observable de la batalla i selecciona una acció per maximitzar la seva esperança de victòria~\cite{russell_norvig,wooldridge_mas}. Una \textbf{política $\boldsymbol{\pi}$} formalitza aquesta presa de decisions assignant una distribució sobre les accions legals a partir de l'observació disponible. Els diferents enfocaments d'Intel·ligència Artificial aborden aquesta política des de perspectives diverses: mitjançant regles heurístiques dissenyades per experts, planificació en línia mitjançant cerca adversària o simulació en arbre (MCTS), o bé aprenentatge de patrons a partir de dades humanes i recompenses per reforç.

Aquest Treball de Final de Màster estudia i compara com cinc paradigmes d'agents afronten aquest problema de decisió sota condicions d'incertesa. L'objectiu no és construir un agent infal·lible, sinó avaluar empíricament implementacions concretes sota pressupostos computacionals explícits, en un entorn reproduïble i amb una escala mostral que permeti quantificar amb rigor la precisió dels resultats.

Per a l'experimentació s'utilitza el simulador \textbf{Pokémon Showdown}~\cite{pokemon_showdown} i la llibreria \textit{poke-env}~\cite{poke_env}, focalitzant la recerca en el format individual \textbf{Random Battle de la Generació~IX} (\texttt{gen9randombattle}). Tot i que la generació procedural d'equips introdueix una forta variabilitat inicial on la composició pot determinar l'avantatge d'una partida, l'execució d'un volum massiu de batalles per enfrontament permet diluir aquest soroll estocàstic i aïllar exclusivament la qualitat de la presa de decisions durant el combat.

Concretament, s'implementen i avaluen cinc formes de decidir sobre el mateix simulador i format:
\begin{enumerate}
    \item \textbf{Heurístiques expertes} (\texttt{v1}--\texttt{v14}): Aquest enfocament es basa en regles de domini dissenyades manualment que assignen una puntuació a cada acció legal segons l'estat del combat. Constitueixen la línia de base determinista i s'estudien mitjançant una seqüència d'ablacions estructurades, on les versions incorporen gradualment càlcul de dany, efectivitat de tipus, perills d'entrada (\emph{hazards}), moviments de preparació (\emph{setup}), gestió d'estats alterats i Teracristal·lització, destacant \texttt{v12} com la variant experta principal.
    \item \textbf{Cerca adversària Minimax d'un torn} (\texttt{v15}--\texttt{v17}): Aquest paradigma construeix una matriu de recompenses d'accions simultànies d'un torn i avalua l'estat resultat amb funcions heurístiques a les fulles. Es comparen tres variants: \texttt{v15} (avaluació basada purament en HP), \texttt{v16} (afegeix bonificacions posicionals per preparació, perills d'entrada i estats) i \texttt{v17} (suma un prior de $+0{,}15$ a l'acció recomanada per \texttt{v14}, reduint la discrepància respecte a l'heurística).
    \item \textbf{Cerca en arbre de Monte Carlo rasa} (\emph{shallow} MCTS \texttt{v18}--\texttt{v20}): Aquest mètode de planificació en línia executa 100 simulacions d'una profunditat rasa de cinc torns mitjançant un simulador local en memòria (\texttt{LocalSim}). Es comparen tres variants: \texttt{v18} (selecció UCB1 i avaluació bàsica d'HP/estats), \texttt{v19} (UCB1 amb fulla millorada per rols i dreceres tàctiques de \texttt{v14}) i \texttt{v20} (selecció \textbf{PUCT} guiada per priors de probabilitat extrets de l'heurística \texttt{v14}).
    \item \textbf{Aprenentatge per imitació} (\texttt{v21}--\texttt{v22}): Aquest enfocament aprèn directament del comportament humà d'alt nivell (Elo~$1800+$) entrenant classificadors XGBoost sobre 1,12 milions de torns. Es comparen dues piles: \texttt{v21} és una versió \textbf{híbrida} (XGBoost decideix la decisió macro d'atacar o canviar, i \texttt{v14} executa l'acció micro concreta) i \texttt{v22} és un model d'\textbf{imitació pura dual} (executa dos caps XGBoost: un per a la decisió macro i un segon model per avaluar atributs de moviments, sense dependre de \texttt{v14}).
    \item \textbf{Aprenentatge per reforç profund} (\emph{MaskablePPO}): Aquest paradigma optimitza automàticament una política de xarxa neuronal mitjançant la interacció i la cerca de recompenses en un entorn amb accions emmascarades. L'entrenament segueix un currículum progressiu que comença contra oponents bàsics (\texttt{random} i \texttt{MaxBasePower}), s'inicialitza per clonatge conductual de l'heurística \texttt{v8} per superar el problema de l'exploració inicial (\emph{cold start}), i finalment s'entrena amb MaskablePPO contra l'heurística experta \texttt{v12}.
\end{enumerate}

Les quatre primeres famílies, juntament amb sis agents de referència, formen una matriu dirigida de \textbf{28 agents}. El PPO no forma part d'aquesta matriu principal: s'avalua en un experiment separat contra \texttt{v12}, l'agent amb millor resultat observat a la lliga. A més, models de llenguatge com \textit{PokéChamp} i \textit{PokéLLMon} s'integren de manera exploratòria.

\section{Objectius del treball}
\label{sec:objectius}

L'objectiu principal d'aquest treball és avaluar i comparar empíricament el rendiment relatiu de les heurístiques expertes, la cerca adversària d'un torn, l'MCTS ras, l'aprenentatge per imitació i PPO en batalles \texttt{gen9randombattle}, sota pressupostos computacionals explícits i un protocol experimental rigorós.

Aquest objectiu general es concreta en els següents objectius específics:
\begin{enumerate}
    \item \textbf{Formalització teòrica del domini:} Formalitzar el combat Pokémon com un procés de decisió estocàstic d'observació parcial (POMDP/POSG) i definir un marc teòric d'avaluació amb incertesa quantificada.
    \item \textbf{Disseny metodològic, infraestructura i telemetria:} Definir el protocol experimental i el model de rànquing Bradley-Terry, construir una infraestructura d'execució paral·lela, escalable i recuperable davant de fallades sobre Pokémon Showdown i \textit{poke-env}, i instrumentar un marc de telemetria per torn per quantificar mètriques explicatives del comportament (predicció de canvis, control del ritme o \emph{tempo}, gestió del risc i ús de la Teracristal·lització).
    \item \textbf{Paradigmes heurístics i de cerca en línia:} Desenvolupar la seqüència d'agents heurístics (\texttt{v1}--\texttt{v14}) com una ablació estructurada de coneixement de domini, i implementar i instrumentar variants de cerca adversària Minimax d'un torn (\texttt{v15}--\texttt{v17}) i MCTS ras de 5 torns (\texttt{v18}--\texttt{v20}) guiats per priors heurístics per analitzar si la cerca millora l'heurística de partida.
    \item \textbf{Agents apresos i integració de LLMs:} Entrenar models d'aprenentatge per imitació XGBoost (\texttt{v21}--\texttt{v22}) sobre repeticions humanes d'alt nivell (Elo~$1800+$) per a decisions macro i micro, optimitzar MaskablePPO amb inicialització per clonatge conductual sobre \texttt{v8} per avaluar si apropa el rendiment a \texttt{v12}, i integrar exploratòriament agents basats en LLM (\textit{PokéChamp} i \textit{PokéLLMon}) analitzant-ne la viabilitat i el cost d'inferència en temps real.
    \item \textbf{Avaluació empírica massiva i rànquings:} Executar i analitzar la matriu competitiva de 28 agents (més de 6,4 milions de batalles), avaluant-ne els rànquings Bradley-Terry i creuant els resultats amb la telemetria per torn per comparar el rendiment relatiu i els patrons de decisió de cada paradigma.
    \item \textbf{Avaluació externa al ladder públic:} Desplegar l'agent heurístic principal (\texttt{v14}) a la classificació pública oficial (\emph{ladder}) de Pokémon Showdown per obtenir una mostra externa de rendiment contra jugadors humans reals en un entorn no controlat.
\end{enumerate}

\section{Estructura de la memòria}

La resta d'aquesta memòria s'organitza en sis capítols interconnectats que desenvolupen el treball des dels fonaments teòrics fins a la discussió dels resultats empírics:

\begin{itemize}
    \item El \textbf{Capítol~\ref{chap:marc_teoric} (Marc teòric i contextualització)} descriu les mecàniques de combat de Pokémon i el format \texttt{gen9randombattle}, formalitza el problema des de la Teoria de Jocs com un procés de decisió estocàstic d'observació parcial (POMDP/POSG), revisa l'estat de l'art en IA aplicada a jocs complexos i detalla els fonaments teòrics dels cinc paradigmes d'agents estudiats.
    \item El \textbf{Capítol~\ref{chap:metodologia} (Metodologia i disseny experimental)} estableix la separació entre el disseny conceptual i la implementació, defineix el gestor de batalles, el protocol d'avaluació de la matriu de 28 agents (més de 6,4M de batalles), el model de rànquing Bradley-Terry, el marc de telemetria per torn per a l'explicabilitat conductual i el control d'amenaces a la validesa.
    \item El \textbf{Capítol~\ref{chap:infraestructura} (Entorn de desenvolupament i infraestructura de còmput)} documenta l'organització del projecte mitjançant un tauler Kanban, l'evolució de l'arquitectura de maquinari (de prototip a l'estació de treball AMD Ryzen 7 5700X3D + RTX 2080), l'orquestració de múltiples servidors Pokémon Showdown en paral·lel, la gestió de persistència, el sistema de notificacions i l'anàlisi de cost econòmic i recursos computacionals.
    \item El \textbf{Capítol~\ref{chap:implementacio} (Implementació dels cinc paradigmes)} descriu el mapa del programari (\texttt{src/}) i tradueix a codi Python el coordinador de batalles, la seqüència d'ablacions heurístiques (\texttt{v1}--\texttt{v14}), la cerca Minimax (\texttt{v15}--\texttt{v17}), l'MCTS ras amb simulador local (\texttt{v18}--\texttt{v20}), els models d'imitació XGBoost (\texttt{v21}--\texttt{v22}), la canalització MaskablePPO amb clonatge conductual, la integració exploratòria de LLMs i l'arquitectura de desplegament operatiu (\emph{Deployment}) per al \emph{ladder} públic.
    \item El \textbf{Capítol~\ref{chap:resultats} (Resultats)} analitza empíricament l'evolució de l'escala heurística, el rendiment de la cerca adversària, MCTS i la imitació, la comparació global dels cinc paradigmes, l'experiment DRL amb PPO contra \texttt{v12}, la validació de \texttt{v14} contra jugadors humans al \emph{ladder} públic i la discussió dels patrons residuals.
    \item El \textbf{Capítol~\ref{chap:conclusions} (Conclusions i línies futures)} respon a la pregunta de recerca principal, sintetitza les sis troballes empíriques clau, exposa les contribucions, analitza les limitacions observades, desenvolupa les línies de treball futur i avalua les consideracions ètiques, socials i d'impacte ambiental.
\end{itemize}
